% Lösungen Einheit 1 von 5 – Bernoulli-Experiment und Bernoulli-Kette (zusammenbau v0.1)
\einheitenkopf{Einheit 1 von 5 · Bernoulli-Experiment und Bernoulli-Kette}

\erg{11}{a) Bernoulli-Kette – zwei Ausgänge, feste Anzahl, gleiches p, unabhängig \quad b) $p = \frac{1}{2} = 0{,}5$ \quad c) $X \sim B(12;\, 0{,}4)$, also $n = 12$ und $p = 0{,}4$ \quad d) $2 \cdot 0{,}15 \cdot 0{,}85 = 0{,}255$ (Pfade GN und NG) \quad e) $3 \cdot \left(\frac{3}{4}\right)^2 \cdot \frac{1}{4} + \left(\frac{3}{4}\right)^3 = \frac{27}{64} + \frac{27}{64} = \frac{27}{32}$ \quad f) Zwei Ausgänge je Brief (beschädigt oder nicht); feste Anzahl der ausgewählten Briefe; gleiche Wahrscheinlichkeit für jeden Brief, weil aus sehr vielen Briefen ausgewählt wird und das Nicht-Zurücklegen kaum etwas ändert; die Briefe sind voneinander unabhängig. \quad g) Nicht binomialverteilt: ohne Zurücklegen aus nur dreißig Laptops ändert sich die Wahrscheinlichkeit für „defekt“ von Zug zu Zug, bei acht von dreißig zu stark, um es zu vernachlässigen; außerdem können höchstens fünf defekte gezogen werden, der Wertebereich reicht nicht bis acht. \quad h) Erste Aussage richtig: eine Drehung hat nur die zwei Ergebnisse A und B. Zweite Aussage falsch: ein Spiel hat drei Ausgänge – Hauptpreis, Trostpreis oder kein Preis (AB oder BA).}
\erg{12}{a) Fehler: Mia prüft nur die zwei Ausgänge. Jan zieht ohne Zurücklegen aus nur fünfzehn Äpfeln, die Wahrscheinlichkeit für „faul“ ändert sich von Zug zu Zug – kein Binomialmodell. \quad b) Weil die Gesamtheit riesig ist, ändert das Herausnehmen von tausend Personen den Anteil kaum; die Wahrscheinlichkeit bleibt praktisch gleich, als würde mit Zurücklegen gezogen. \quad c) $1 - 0{,}001^2 = 0{,}999999$ \quad d) Baum mit zwei Stufen, an jedem Ast G mit $0{,}25$ und V mit $0{,}75$}

12d)

\baumzwei{G/0{,}25,V/0{,}75}{G/0{,}25,V/0{,}75}{G/0{,}25,V/0{,}75}

